% !TeX root = ../main.tex

\newlecture[Fourier Series: Formulation]

\subsection{2.1 Introduction}
In order to solve a PDE, we must be able to represent its solution, which is a function in space (and time). We could attempt to describe the function in terms of the value that it takes at each point in space (and time), but this representation is not convenient for solving or analyzing PDEs. A much more convenient approach is to represent the function as a linear combination of known functions, so that the function can be identified by a set of coefficients. Namely, we represent the solution $u : \Omega \to \R$ as $u(x) = \sum_{n=1}^{\infty} \hat{u}_n \phi_n(x)$, where $\phi_n$, $n = 1, 2, \dots$, are the known functions and $\hat{u}_n$, $n = 1, 2, \dots$, are the coefficients to be determined. One class of infinite series that is particularly useful for representing solutions of PDEs is \textit{Fourier series}. In this lecture, we introduce Fourier series so that we can use them to solve PDEs in the following lectures. In this lecture, we focus on the construction of Fourier series; we will analyze their theoretical properties in the next lecture.

\subsection{2.2 Periodic, even, and odd functions}
As periodic functions play an important role in the construction of Fourier series, we first provide a few definitions related to the periodicity of functions.

\begin{definition}[Periodic Function][2.1]
    A function $f : \R \to \R$ is said to be \textbf{$z$-periodic} if
    \[f(x + z) = f(x) \quad \forall x \in \R.\]
    The real number $z$ is called a \textit{period}. The smallest positive real number $z$ such that the relationship holds is called the \textit{fundamental period}.
\end{definition}

\begin{example}[Periodicity of Sine Functions][2.2]
    Perhaps one of the most common periodic functions is the sine function. The function $f(x) = \sin(x)$ is $2\pi$-periodic because
    \[\sin(x + 2\pi) = \sin(x) \quad \forall x \in \R.\]
    Its fundamental period is $2\pi$, because the relationship does not hold for any shorter period. Note that any integer multiple of the fundamental period is also a period of the function; e.g., $z = 2\pi, 4\pi, 6\pi$.
\end{example}

\begin{example}[Periodicity of a Sum of Sinusoidal Functions][2.3]
    Consider $f(x) = \cos(x) + \sin(2x)$. The functions $\cos(x)$ and $\sin(2x)$ are periodic functions with fundamental periods $2\pi$ and $\pi$, respectively. The fundamental period of the sum of the functions is $2\pi$, the least common period of the two functions.
\end{example}

We also recall the definition of even and odd functions.

\begin{definition}[Even Function][2.4]
    A function $f : (-a, a) \to \R$ is an \textbf{even} function if for all $x \in (-a, a)$,
    \[f(-x) = f(x).\]
\end{definition}

\begin{definition}[Odd Function][2.5]
    A function $f : (-a, a) \to \R$ is an \textbf{odd} function if for all $x \in (-a, a)$,
    \[f(-x) = -f(x).\]
\end{definition}

If $f$ is even, then the graph $y = f(x)$ is symmetric about the $y$-axis; the left-hand side is the mirror image of the right-hand side. If $f$ is odd, then the graph $y = f(x)$ is symmetric about the origin. We provide a few examples.

\begin{example}[][2.6]
    For any $\lambda \geq 0$, the function $f(x) = \cos(\lambda x)$ is an even function because $f(-x) = \cos(-\lambda x) = \cos(\lambda x) = f(x)$. For any $\lambda \geq 0$, the function $f(x) = \sin(\lambda x)$ is an odd function because $f(-x) = \sin(-\lambda x) = -\sin(\lambda x) = -f(x)$.
\end{example}

\begin{example}[][2.7]
    A monomial $x^k$ is an even function if $k$ is even and an odd function if $k$ is odd.
\end{example}

We also note a few properties of even and odd functions under several operations; each result can be readily proven from the definition of even and odd functions:
\begin{itemize}
    \item \textit{Summation.} The sum of two even functions is even. The sum of two odd functions is odd. However, the sum of an even function and an odd function is in general neither even nor odd.
    \item \textit{Multiplication.} The product of two even functions or two odd functions is even. The product of an even function and an odd function is odd.
    \item \textit{Differentiation.} The derivative of an odd function is even, and the derivative of an even function is odd. The result follows from the chain rule.
    \item \textit{Integration.} If $f$ is an even function, then the integral $g(x) = \int_{\xi=0}^{x} f(\xi) d\xi$ is odd. If $f$ is an odd function, then the integral $g(x) = \int_{\xi=0}^{x} f(\xi) d\xi$ is even.
    \item \textit{Integrals.} If $f$ is odd, then $\int_{-a}^{a} f(x) dx = 0$. If $f$ is even, then $\int_{-a}^{a} f(x) dx = 2 \int_0^a f(x) dx$.
    \item \textit{Behavior at origin.} If $f$ is even and is differentiable at $x = 0$, then its derivative must vanish: i.e., $f'(x = 0) = 0$. If $f$ is odd and is continuous at $x = 0$, then its value must vanish: i.e., $f(x = 0) = 0$.
\end{itemize}

Given a function defined on a finite interval, we may also ``extend'' the function to the entire real line $\R$ to construct a periodic function on $\R$. We introduce a few different extensions.

\begin{definition}[Periodic Extension][2.8]
    Given a function $f : (0, a) \to \R$, its \textbf{periodic extension} to $\R$ is the function $f^{\mathrm{p}} : \R \to \R$ such that
    \[f^{\mathrm{p}}(x) = f(x), \quad x \in (0, a),\]
    and has a period $a$ (i.e., $f^{\mathrm{p}}(x + a) = f^{\mathrm{p}}(x)$, $\forall x \in \R$, or equivalently $f^{\mathrm{p}}(x + na) = f(x)$, $\forall x \in (0, a)$ and $\forall n \in \Z$).
\end{definition}

\begin{definition}[Even Periodic Extension][2.9]
    Given a function $f : (0, a) \to \R$, its \textbf{even periodic extension} is the function $f^{\mathrm{e,p}} : \R \to \R$ such that
    \begin{align*}
        f^{\mathrm{e,p}}(x) &= f(x), \quad x \in (0, a), \\
        f^{\mathrm{e,p}}(x) &= f(-x), \quad x \in (-a, 0),
    \end{align*}
    and has a period of $2a$ (i.e., $f^{\mathrm{e,p}}(x + 2a) = f^{\mathrm{e,p}}(x)$, $\forall x \in \R$).
\end{definition}

\begin{definition}[Odd Periodic Extension][2.10]
    Given a function $f : (0, a) \to \R$, its \textbf{odd periodic extension} is the function $f^{\mathrm{o,p}} : \R \to \R$ such that
    \begin{align*}
        f^{\mathrm{o,p}}(x) &= f(x), \quad x \in (0, a), \\
        f^{\mathrm{o,p}}(x) &= -f(-x), \quad x \in (-a, 0),
    \end{align*}
    and has a period of $2a$ (i.e., $f^{\mathrm{o,p}}(x + 2a) = f^{\mathrm{o,p}}(x)$, $\forall x \in \R$).
\end{definition}

In the above definitions of periodic, even-periodic, and odd-periodic extensions, we do not specify the values of the extensions at the endpoints $x = na$, $n \in \Z$. As we will see shortly, the lack of specification does not cause a problem in our discussion of Fourier series.

We demonstrate these three periodic extensions using a concrete example.

\begin{example}[Periodic Extensions][2.11]
    Let $f(x) = x$, $x \in (0, 1)$. Figure 2.1 shows the periodic extension, even periodic extension, and odd periodic extension for the function. To construct the periodic extension, we simply repeat the function, with a period of $1$. We construct the even periodic extension in two steps: we first mirror the function about the $y$-axis to construct the even extension; we then construct the periodic extension of the even extension with a period of $2$. By definition, the even (or odd) extension of a function is defined on a symmetric interval (e.g., $[-1, 1]$). We also note that the period of the even periodic extension is twice the domain length of the original function to match the domain length of the even extension. Similarly, we construct the odd periodic extension in two steps: we first mirror the function about the origin to construct the odd extension; we then construct the periodic extension of the odd extension with a period of $2$.
\end{example}

\lecfig{fourier_series}{4}{79 582 79 81}{Figure 2.1: Periodic extensions of $f(x) = x$ for $x \in [0, 1]$.}

\subsection{2.3 Fourier series}
We now introduce Fourier series, which provide a representation of a periodic function as a linear combination of sinusoidal functions. For simplicity, we consider $2$-periodic functions; we will later generalize the formulation to an arbitrary period in Section 2.5.

\begin{definition}[Fourier Series][2.12]
    Let $f : \R \to \R$ be a $2$-periodic function. The \textbf{Fourier series} of the function is
    \begin{equation*}
        f(x) \sim \frac{1}{2} a_0 + \sum_{n=1}^{\infty} \left(a_n \cos(n\pi x) + b_n \sin(n\pi x)\right), \tag{2.1}
    \end{equation*}
    with Fourier coefficients
    \begin{equation*}
        \begin{aligned}
            a_n &= \int_{-1}^{1} \cos(n\pi x) f(x) dx, \quad n = 0, 1, 2, \dots, \\
            b_n &= \int_{-1}^{1} \sin(n\pi x) f(x) dx, \quad n = 1, 2, \dots.
        \end{aligned} \tag{2.2}
    \end{equation*}
\end{definition}

\begin{definition}[Truncated Fourier Series][2.13]
    Let $f : \R \to \R$ be a $2$-periodic function. The \textbf{$N$-term truncated Fourier series} of the function is
    \[s_N(x) \equiv \frac{1}{2} a_0 + \sum_{n=1}^{N} \left(a_n \cos(n\pi x) + b_n \sin(n\pi x)\right),\]
    with coefficients given by (2.2).
\end{definition}

We make some observations about Fourier series.
\begin{enumerate}
    \item Because we have not verified, and in general will not be able to verify, that the Fourier series of $f$ is the same as $f$, we use the notation ``$\sim$'' instead of ``$=$'' in (2.1). In the next lecture, we will answer various questions about the convergence of Fourier series including the following: does the series converge? In what sense does the series converge? How smooth does the function $f$ have to be for the series to converge?
    \item The set of functions
    \[\left\{\frac{1}{\sqrt{2}}, \sin(\pi x), \cos(\pi x), \sin(2\pi x), \cos(2\pi x), \dots\right\}\]
    is \textit{orthonormal} with respect to the $L^2$ inner product, $(w, v) \equiv \int_{-1}^{1} wv\,dx$. For all $n, m \in \N_{>0}$,
    \begin{align*}
        \int_{-1}^{1} \frac{1}{\sqrt{2}} \cos(m\pi x) dx = \int_{-1}^{1} \frac{1}{\sqrt{2}} \sin(m\pi x) dx &= 0, \\
        \int_{-1}^{1} \cos(n\pi x) \cos(m\pi x) dx &= \delta_{mn}, \\
        \int_{-1}^{1} \sin(n\pi x) \sin(m\pi x) dx &= \delta_{mn}, \\
        \int_{-1}^{1} \sin(n\pi x) \cos(m\pi x) dx &= 0,
    \end{align*}
    where $\delta_{mn}$ is the Kronecker delta; i.e., $\delta_{mn} = 1$ if $m = n$ and $\delta_{mn} = 0$ if $m \neq n$.
    \item The expression for the coefficients (2.2) is a consequence of the orthonormality of the sinusoidal functions. For instance, to obtain the expression for $a_2$, we multiply the Fourier series (2.1) by $\cos(2\pi x)$, integrate the expression over $(-1, 1)$, and appeal to the orthonormality: i.e.,
    \begin{align*}
        \int_{-1}^{1} \cos(2\pi x) f(x) dx &= \int_{-1}^{1} \cos(2\pi x) \left(\frac{1}{2} a_0 + \sum_{n=1}^{\infty} \left(a_n \cos(n\pi x) + b_n \sin(n\pi x)\right)\right) dx \\
        &= a_0 \int_{-1}^{1} \frac{1}{2} \cos(2\pi x) dx + \sum_{n=1}^{\infty} a_n \int_{-1}^{1} \cos(2\pi x) \cos(n\pi x) dx \\
        &\qquad + \sum_{n=1}^{\infty} b_n \int_{-1}^{1} \cos(2\pi x) \sin(n\pi x) dx \\
        &= \sum_{n=1}^{\infty} a_n \delta_{2n} = a_2.
    \end{align*}
    All integrals with the exception of the one multiplied by $a_2$ vanish due to the mutual orthogonality of the sine and cosine functions. We obtain the expression for all other $a_n$ and $b_n$ using the same orthogonality argument.
    \item We can also find the Fourier series of a function $f$ defined on $(-1, 1)$ (instead of a $2$-periodic function on $\R$). The Fourier series of $f : (-1, 1) \to \R$ is given by (2.1) and it converges to (i) $f$ on $(-1, 1)$ and (ii) its periodic extension $f^{\mathrm{p}}$ on $\R$.
    \item The truncated Fourier series $s_N$ of $f$ is the $N$-term partial sum of the Fourier series of $f$.
\end{enumerate}

We now consider an example:

\begin{example}[The Fourier Series of $f(x) = x$ on $(-1, 1)$][2.14]
    We consider the $2$-periodic extension of the function $f(x) = x$ for $x \in [-1, 1]$. We wish to find the Fourier series representation of the function. The coefficient $a_n$, $n \in \N_{\geq 0}$, is given by
    \[a_n = \int_{-1}^{1} \cos(n\pi x) f(x) dx = \int_{-1}^{1} \cos(n\pi x) x\,dx = 0;\]
    the integral evaluates to $0$ because the integrand is an odd function, and the domain of integration is symmetric. The Fourier coefficient $b_n$ is given by
    \begin{align*}
        b_n &= \int_{-1}^{1} \sin(n\pi x) f(x) dx = \int_{-1}^{1} \sin(n\pi x) x\,dx = \left[-\frac{1}{n\pi} \cos(n\pi x) x\right]_{x=-1}^{1} + \int_{-1}^{1} \frac{1}{n\pi} \cos(n\pi x) dx \\
        &= -\frac{1}{n\pi} \left(\cos(n\pi) + \cos(-n\pi)\right) + \frac{1}{n^2\pi^2} \sin(n\pi x) \Big|_{x=-1}^{1} = -\frac{2}{n\pi} \cos(n\pi) = -\frac{2}{n\pi} (-1)^n \\
        &= \frac{2}{n\pi} (-1)^{n+1};
    \end{align*}
    here, the third equality follows from integration by parts, and the penultimate equality follows from $\cos(n\pi) = (-1)^n$, $\forall n \in \Z$. It follows that
    \begin{align*}
        f(x) \sim \sum_{n=1}^{\infty} b_n \sin(n\pi x) &= \sum_{n=1}^{\infty} \frac{2(-1)^{n+1}}{n\pi} \sin(n\pi x) \\
        &= \frac{2}{\pi} \left(\sin(\pi x) - \frac{1}{2} \sin(2\pi x) + \frac{1}{3} \sin(3\pi x) - \cdots\right).
    \end{align*}
    The series contains only sine terms (and not cosine terms) because the function $f(x) = x$ is an odd function. Figure 2.2 illustrates the Fourier series approximation of the function. We observe that the weighted sum of the sine functions at least visually converges to $f$. (As noted earlier, we will analyze the convergence of Fourier series in the next lecture.)
\end{example}

\lecfig[0.5\linewidth]{fourier_series}{6}{184 560 209 72}{Figure 2.2: Truncated Fourier series of $f(x) = x$ for $x \in [-1, 1]$.}

\subsection{2.4 Fourier sine and cosine series}
The Fourier series provides a representation for any periodic function. It is often also useful to construct a representation specialized to odd or even periodic functions. A Fourier sine series is the Fourier series associated with odd periodic functions, and a Fourier cosine series is the Fourier series associated with even periodic functions.

\begin{definition}[Fourier Cosine Series][2.15]
    Let $f : \R \to \R$ be an even $2$-periodic function. The \textbf{Fourier cosine series} of $f$ is given by
    \begin{equation*}
        f(x) \sim \frac{1}{2} \hat{f}_0 + \sum_{n=1}^{\infty} \hat{f}_n \cos(n\pi x) \tag{2.3}
    \end{equation*}
    where
    \[\hat{f}_n = \int_{-1}^{1} \cos(n\pi x) f(x) dx = 2 \int_0^1 \cos(n\pi x) f(x) dx, \quad n \in \N_{\geq 0}.\]
\end{definition}

\begin{definition}[Fourier Sine Series][2.16]
    Let $f : \R \to \R$ be an odd $2$-periodic function. The \textbf{Fourier sine series} of $f$ is given by
    \begin{equation*}
        f(x) \sim \sum_{n=1}^{\infty} \hat{f}_n \sin(n\pi x) \tag{2.4}
    \end{equation*}
    where
    \[\hat{f}_n = \int_{-1}^{1} \sin(n\pi x) f(x) dx = 2 \int_0^1 \sin(n\pi x) f(x) dx, \quad n \in \N_{>0}.\]
\end{definition}

We make some observations:
\begin{enumerate}
    \item The domain of integration for the coefficients can be reduced from $[-1, 1]$ to $[0, 1]$ due to the symmetry of even functions for the Fourier cosine series (and due to the antisymmetry of odd functions for the Fourier sine series).
    \item Similar to the (full) Fourier series, we can also find the Fourier cosine series of a function $f$ defined on $(0, 1)$ (instead of an even $2$-periodic function on the entire $\R$). The Fourier cosine series of $f : (0, 1) \to \R$ is given by (2.3) and it converges to (i) $f$ on $(0, 1)$ and (ii) its even periodic extension $f^{\mathrm{e,p}}$ on the entire $\R$.
    \item Similarly, we can find the Fourier sine series of a function $f$ defined on $(0, 1)$ (instead of an odd $2$-periodic function on the entire $\R$). The Fourier sine series of $f : (0, 1) \to \R$ is given by (2.4) and it converges to (i) $f$ on $(0, 1)$ and (ii) its odd periodic extension $f^{\mathrm{o,p}}$ on the entire $\R$.
    \item We use the notation $(\hat{f}_n)_{n=0}^{\infty}$ to denote the Fourier sine/cosine coefficients of $f$. The notation makes the connection between the function and its Fourier coefficients more explicit, which will be helpful when we work with Fourier series associated with multiple functions. (Unfortunately, this notation cannot be used for the (full) Fourier series, as we require two sets of coefficients for sine and cosine.)
\end{enumerate}

We now consider an example of Fourier cosine and sine series.

\begin{example}[Fourier Cosine Series of $f(x) = x$ on $(0, 1)$][2.17]
    We consider the Fourier cosine series representation of $f(x) = x$ for $x \in (0, 1)$. The coefficient $\hat{f}_0$ of the Fourier cosine series is given by
    \[\hat{f}_0 = 2 \int_0^1 x\,dx = 1.\]
    For $n > 0$, we obtain
    \begin{align*}
        \hat{f}_n &= 2 \int_0^1 x \cos(n\pi x) dx = 2 \left[\frac{\cos(n\pi x)}{n^2\pi^2} + \frac{x \sin(n\pi x)}{n\pi}\right]_{x=0}^{1} = \frac{2}{n^2\pi^2} (\cos(n\pi) - 1) \\
        &= \frac{2}{n^2\pi^2} ((-1)^n - 1) = \begin{cases} -\frac{4}{n^2\pi^2} & n \in \text{odd}, \\ 0 & n \in \text{even}. \end{cases}
    \end{align*}
    The Fourier cosine series of $f$ is therefore
    \begin{align*}
        f(x) \sim \frac{1}{2} \hat{f}_0 + \sum_{n=1}^{\infty} \hat{f}_n \cos(n\pi x) &= \frac{1}{2} + \sum_{k=1}^{\infty} \hat{f}_{2k-1} \cos((2k-1)\pi x) \\
        &= \frac{1}{2} - \sum_{k=1}^{\infty} \frac{4}{(2k-1)^2\pi^2} \cos((2k-1)\pi x) \\
        &= \frac{1}{2} - \frac{4}{\pi^2} \left(\cos(\pi x) + \frac{1}{9} \cos(3\pi x) + \frac{1}{25} \cos(5\pi x) + \cdots\right).
    \end{align*}
    Figure 2.3(a) illustrates the Fourier cosine series representation of the function. We observe that the series (at least visually) approximates the function $f(x) = x$ on $(0, 1)$. Note that the Fourier cosine series approximates the even periodic extension of $f$.
\end{example}

\lecfig{fourier_series}{8}{115 544 101 83}{Figure 2.3: Truncated Fourier cosine and sine series of $f(x) = x$ for $x \in [0, 1]$.}

\begin{example}[Fourier Sine Series Approximation of $f(x) = x$ on $(0, 1)$][2.18]
    We consider the Fourier sine series representation of $f(x) = x$ for $x \in (0, 1)$. The coefficients of the Fourier sine series are given by
    \[\hat{f}_n = 2 \int_0^1 x \sin(n\pi x) dx = 2 \left[\frac{\sin(n\pi x)}{n^2\pi^2} - \frac{x \cos(n\pi x)}{n\pi}\right]_{x=0}^{1} = -\frac{2}{n\pi} \cos(n\pi) = \frac{2}{n\pi} (-1)^{n+1}.\]
    The Fourier sine series of $f$ is therefore
    \[f(x) \sim \sum_{n=1}^{\infty} \hat{f}_n \sin(n\pi x) = \sum_{n=1}^{\infty} \frac{2(-1)^{n+1}}{n\pi} \sin(n\pi x) = \frac{2}{\pi} \left(\sin(\pi x) - \frac{1}{2} \sin(2\pi x) + \frac{1}{3} \sin(3\pi x) - \cdots\right).\]
    Not surprisingly, the Fourier sine series of $f(x) = x$ on $(0, 1)$ considered here is the same as the Fourier series of $g(x) = x$ on $(-1, 1)$ considered in Example 2.14 because the odd periodic extension of $f$ is the same as the periodic extension of $g$. Figure 2.3(b) illustrates the Fourier sine series representation of the function. Note that the Fourier sine series approximates the odd periodic extension of $f$.
\end{example}

\subsection{2.5 Fourier series on arbitrary intervals}
We now consider Fourier series for periodic functions of arbitrary periods. Given a $2L$-periodic function $f : \R \to \R$, the Fourier series of $f$ results from simple adjustments of the arguments of the sine and cosine functions to the period of $2L$:
\[f(x) \sim \frac{1}{2} a_0 + \sum_{n=1}^{\infty} \left(a_n \cos\left(\frac{n\pi x}{L}\right) + b_n \sin\left(\frac{n\pi x}{L}\right)\right),\]
with coefficients
\begin{align*}
    a_n &= \frac{1}{L} \int_{-L}^{L} \cos\left(\frac{n\pi x}{L}\right) f(x) dx, \quad n = 0, 1, 2, \dots, \\
    b_n &= \frac{1}{L} \int_{-L}^{L} \sin\left(\frac{n\pi x}{L}\right) f(x) dx, \quad n = 1, 2, \dots.
\end{align*}
The normalization factor of the coefficients for the Fourier series on $(-L, L)$ is scaled by $1/L$ relative to the Fourier series on $(-1, 1)$ because the orthogonality relations are
\begin{align*}
    \frac{1}{L} \int_{-L}^{L} \sin\left(\frac{n\pi x}{L}\right) \cos\left(\frac{m\pi x}{L}\right) dx &= 0, \\
    \frac{1}{L} \int_{-L}^{L} \sin\left(\frac{n\pi x}{L}\right) \sin\left(\frac{m\pi x}{L}\right) dx &= \delta_{mn}, \\
    \frac{1}{L} \int_{-L}^{L} \cos\left(\frac{n\pi x}{L}\right) \cos\left(\frac{m\pi x}{L}\right) dx &= \delta_{mn},
\end{align*}
for all $m, n \in \N_{>0}$. The Fourier series of a function $f : (-L, L) \to \R$ is also given by the same expression, and the series converges to the $2L$-periodic extension of $f$ on the entire $\R$.

Similarly, the Fourier cosine series of an even $2L$-periodic function $f : \R \to \R$ is given by
\[f(x) \sim \frac{1}{2} \hat{f}_0 + \sum_{n=1}^{\infty} \hat{f}_n \cos\left(\frac{n\pi x}{L}\right)\]
with Fourier cosine coefficients
\[\hat{f}_n = \frac{2}{L} \int_0^L \cos\left(\frac{n\pi x}{L}\right) f(x) dx, \quad n = 0, 1, 2, \dots.\]
The Fourier cosine series of a function $f : (0, L) \to \R$ is also given by the same expression, and the series represents the even $2L$-periodic extension of $f$ on the entire $\R$.

The Fourier sine series of an odd $2L$-periodic function $f : \R \to \R$ is given by
\[f(x) \sim \sum_{n=1}^{\infty} \hat{f}_n \sin\left(\frac{n\pi x}{L}\right)\]
with Fourier sine coefficients
\[\hat{f}_n = \frac{2}{L} \int_0^L \sin\left(\frac{n\pi x}{L}\right) f(x) dx, \quad n = 1, 2, \dots.\]
The Fourier sine series of a function $f : (0, L) \to \R$ is also given by the same expression, and the series converges to the odd $2L$-periodic extension of $f$ on the entire $\R$.

\subsection{2.6 Summary}
We summarize key points of this lecture:
\begin{enumerate}
    \item Fourier series provide representations of periodic functions $f : \R \to \R$ in terms of sinusoidal functions. To find the coefficients of the series, we multiply $f$ by the sinusoidal functions and integrate over the period; this is a consequence of the orthonormality of the sinusoidal functions.
    \item Fourier cosine and sine series provide representations of even and odd periodic functions, respectively.
    \item Given $f : (-L, L) \to \R$, its (full) Fourier series approximates the periodic extension of $f$. Given $f : (0, L) \to \R$, its Fourier cosine and sine series approximate the even periodic and odd periodic extensions, respectively, of $f$.
    \item Fourier series can be applied to periodic functions of an arbitrary period by adjusting the period of the sinusoidal functions and the associated coefficients accordingly.
\end{enumerate}

\subsection{2.A Complex form of Fourier series}
The Fourier series can be written in a more compact form using complex exponentials. To begin, we recall the de Moivre formulas
\[\sin(x) = \frac{1}{2i} (\exp(ix) - \exp(-ix)) \quad \text{and} \quad \cos(x) = \frac{1}{2} (\exp(ix) + \exp(-ix)).\]
It follows that the set of sinusoidal functions
\[\{1, \cos(\pi x), \sin(\pi x), \cos(2\pi x), \sin(2\pi x), \dots\}\]
can be replaced by the complex exponentials
\[\{1, \exp(-i\pi x), \exp(i\pi x), \exp(-i2\pi x), \exp(i2\pi x), \dots\} = \{\exp(in\pi x)\}_{n=-\infty}^{\infty}.\]
The complex form of the Fourier series is as follows:

\begin{definition}[Complex Form of Fourier Series][2.19]
    Let $f : \R \to \R$ be a $2$-periodic function. The \textbf{complex form of the Fourier series} of $f$ is
    \begin{equation*}
        f(x) \sim \sum_{n=-\infty}^{\infty} \hat{f}_n \exp(in\pi x), \tag{2.5}
    \end{equation*}
    with coefficients
    \[\hat{f}_n = \frac{1}{2} \int_{-1}^{1} \exp(-in\pi x) f(x) dx.\]
\end{definition}

As with the real Fourier series of $f$, the expression for the coefficients follows from the orthogonality relationship. Specifically, the orthogonality relationship for the complex exponentials $\exp(-im\pi x)$ and $\exp(in\pi x)$ for $m \neq n$ is
\begin{align*}
    \int_{-1}^{1} \exp(-im\pi x) \exp(in\pi x) dx &= \int_{-1}^{1} \exp(i(n - m)\pi x) dx \\
    &= \frac{1}{i\pi(n - m)} \left(\exp(i(n - m)\pi) - \exp(-i(n - m)\pi)\right) \\
    &= \frac{2}{\pi(n - m)} \sin((n - m)\pi) = 0,
\end{align*}
and for $m = n$ is
\[\int_{-1}^{1} \exp(i(n - m)\pi x) dx = \int_{-1}^{1} dx = 2.\]
The complex form of Fourier series (2.5) immediately follows from the orthogonality relationship.

For a $2L$-periodic function $f : \R \to \R$, the complex form of the Fourier series is given by
\[f(x) \sim \sum_{n=-\infty}^{\infty} \hat{f}_n \exp\left(\frac{in\pi x}{L}\right),\]
with complex Fourier coefficients
\[\hat{f}_n = \frac{1}{2L} \int_{-L}^{L} \exp\left(\frac{-in\pi x}{L}\right) f(x) dx.\]
The Fourier series of a function $f : (-L, L) \to \R$ is also given by the same expression, and the series converges to the $2L$-periodic extension of $f$ on the entire $\R$.

The complex form of Fourier series is more concise and arguably more elegant than the real form of Fourier series, which requires separate coefficients for the sine and cosine modes. However, in this course we mostly use the real form of Fourier series; this is mainly because Fourier sine and cosine series are far more frequently encountered than the (full) Fourier series in the analysis of PDEs, and the sine/cosine series also admit a compact representation.

\subsection{2.B List of useful integrals}
We list a few indefinite integrals that are useful for the evaluation of Fourier series.
\begin{enumerate}
    \item $\displaystyle\int \sin(\alpha x) dx = -\frac{\cos(\alpha x)}{\alpha}$
    \item $\displaystyle\int \cos(\alpha x) dx = \frac{\sin(\alpha x)}{\alpha}$
    \item $\displaystyle\int x \sin(\alpha x) dx = \frac{\sin(\alpha x)}{\alpha^2} - \frac{x \cos(\alpha x)}{\alpha}$
    \item $\displaystyle\int x \cos(\alpha x) dx = \frac{\cos(\alpha x)}{\alpha^2} + \frac{x \sin(\alpha x)}{\alpha}$
    \item $\displaystyle\int x^2 \sin(\alpha x) dx = \frac{2x \sin(\alpha x)}{\alpha^2} + \frac{(2 - \alpha^2 x^2) \cos(\alpha x)}{\alpha^3}$
    \item $\displaystyle\int x^2 \cos(\alpha x) dx = \frac{2x \cos(\alpha x)}{\alpha^2} + \frac{(\alpha^2 x^2 - 2) \sin(\alpha x)}{\alpha^3}$
    \item $\displaystyle\int \sin(\alpha x) \sin(\beta x) dx = \frac{\sin((\alpha - \beta)x)}{2(\alpha - \beta)} - \frac{\sin((\alpha + \beta)x)}{2(\alpha + \beta)}, \quad \alpha \neq \beta$
    \item $\displaystyle\int \cos(\alpha x) \sin(\beta x) dx = \frac{\cos((\alpha - \beta)x)}{2(\alpha - \beta)} - \frac{\cos((\alpha + \beta)x)}{2(\alpha + \beta)}, \quad \alpha \neq \beta$
    \item $\displaystyle\int \cos(\alpha x) \cos(\beta x) dx = \frac{\sin((\alpha - \beta)x)}{2(\alpha - \beta)} + \frac{\sin((\alpha + \beta)x)}{2(\alpha + \beta)}, \quad \alpha \neq \beta$
    \item $\displaystyle\int \sin^2(\alpha x) dx = \frac{x}{2} - \frac{\sin(2\alpha x)}{4\alpha}$
    \item $\displaystyle\int \cos(\alpha x) \sin(\alpha x) dx = \frac{\sin^2(\alpha x)}{2\alpha}$
    \item $\displaystyle\int \cos^2(\alpha x) dx = \frac{x}{2} + \frac{\sin(2\alpha x)}{4\alpha}$
    \item $\cos(n\pi) = (-1)^n$, $n \in \Z$
    \item $\sin(n\pi) = 0$, $n \in \Z$
\end{enumerate}
